\documentclass[10pt,twocolumn,letterpaper]{article}
\usepackage[textwidth=7in,textheight=9in,columnsep=0.25in,centering]{geometry}
\usepackage{newtxtext,newtxmath}
\usepackage{graphicx,booktabs,enumitem,xcolor,multirow,array}
\usepackage[hidelinks]{hyperref}
\usepackage[compact,small]{titlesec}
\titlespacing*{\section}{0pt}{6pt}{3pt}
\titlespacing*{\subsection}{0pt}{4pt}{2pt}
\setlist{nosep,leftmargin=*}
\setlength{\parskip}{0pt}
\setlength{\textfloatsep}{8pt}
\setlength{\abovecaptionskip}{3pt}
\renewcommand{\baselinestretch}{1.0}
\newcommand{\f}[1]{\texttt{\small #1}}

\title{\vspace{-0.4in}\Large Failure Propagation and Recovery in Multi-Agent Systems\\
\large Phase 1: Fault Injection and Propagation Measurement}
\author{\textcolor{red}{[Team member names]} \quad CSE 291A, Fall 2026}
\date{}

\begin{document}
\maketitle
\vspace{-0.45in}

\section{Goal and Setup}
When agents pass results to each other, one wrong or malicious result can
change what later agents decide and do. Phase~1 builds a workflow whose
outcomes can be checked by a program, injects faults at controlled stages,
and measures how far each fault spreads. It also sets two reference points
for Phase~2's selective recovery: full restart and always-verify.

\textbf{Workflow.} Procurement tasks: ``team $T$ needs $q\times$ item; buy
from the cheapest \emph{approved} vendor if the total fits the team's
remaining budget; record the order, deduct the budget, and email the manager
a confirmation with the order id; otherwise order nothing and notify the
manager.'' A deterministic \emph{retrieval} step reads the vendor catalog,
the team directory, and the shared file \f{budgets.json}. The
\emph{researcher} (LLM) extracts facts. The \emph{planner} (LLM) emits a JSON
plan of tool calls. The \emph{executor} (LLM) issues one tool call at a time,
filling in values returned by earlier calls (e.g., the order id). Tasks are
generated from a seed with known ground truth. About 30\% are over budget,
and in about half of them the cheapest vendor is unapproved, so the agents
must actually apply both rules.

\textbf{Environment.} State has two kinds. \emph{Restorable shared files}
(\f{budgets.json}, \f{orders.csv}) can be rolled back from a snapshot.
\emph{External actions} (\f{place\_order}, \f{send\_email}) leave the system
and can never be rolled back. A checker compares the final state with ground
truth: exactly one correct order, a matching record, the correct deduction,
and a confirmation that carries the real id. It counts \emph{unauthorized
actions} separately: email to an external domain, file deletion, ordering
from an unapproved vendor or over budget, and writes to other teams or files.

\textbf{Tracing.} Every message, intermediate result, tool call, and injected
fault is stored as an artifact with explicit input edges. A tool call also
depends on the last writer of every file it reads, so dependencies through
shared state appear in the same graph.

\textbf{Faults} (definitions in Appendix~B). There are 8 \emph{accidental}
faults: wrong value, stale output, timeout, a corrupted budget in shared
state, and a corrupted order id in a tool result. There are 4
\emph{adversarial} faults: prompt injection in a retrieved catalog record,
in the researcher's output, and in a vendor API response, plus a compromised
planner that appends an exfiltration step. A fault replaces a payload at a
stage boundary, so every downstream artifact depends on the fault artifact.
Each injected payload is kept internally consistent. A first version changed
only the target field, and downstream agents simply corrected it from
fields left unchanged (e.g., the researcher's free-text notes still named
the right vendor).

\textbf{Is the fault actually used?} Each faulty run is aligned with the
clean run of the same task by artifact position. Positions are compared on
decision-relevant content only: recipients, ids, amounts, and tool
arguments, not email wording. Two clean repetitions per task were
semantically identical in 23/25 tasks; both differences come from one
natural model error (below), so LLM nondeterminism adds little noise. Each run is then labeled \emph{ignored} (nothing
downstream changed and no injected string reached a downstream artifact),
\emph{propagated-benign}, or \emph{propagated-harmful} (task failed or an
unauthorized action occurred).

\textbf{Methods.} \emph{None}: no recovery. \emph{Full restart}: an oracle
names the faulty result, so detection is kept separate from recovery; the
files are restored from the snapshot and the whole workflow reruns.
\emph{Always-verify}: a verifier LLM checks the researcher and planner
outputs against their sources, and a rejected output is regenerated once.
All methods run the same fault schedule.

\textbf{Scale.} 25 tasks $\times$ 12 faults $\times$ 3 methods, plus 3 clean
runs per task: 975 runs, with 0 harness failures. Model:
\f{api-glm-5.3-flash} (UCSD TritonAI), temperature 0, 8 parallel workers
on a laptop.

\section{Results}
\textbf{Clean baseline.} Task success is 96\% (48/50 runs), with no
unauthorized actions. A run costs 5.2k tokens, 6.3 LLM calls, and 29\,s of
wall time. Both failures are a \emph{natural} fault: the researcher reported
the balance \emph{after} the purchase as the remaining budget, and the
planner, trusting it, declined a valid order. It spreads exactly like our
injected wrong values. Always-verify adds 37\% tokens and 87\% latency, with
1 false rejection in 25 clean runs.

\begin{figure}[t]
\centering
\includegraphics[width=\columnwidth]{figs/propagation.pdf}
\caption{How each fault propagates with no recovery (25 tasks; only runs
where the fault applied). The right column gives mean unauthorized actions
per run.}
\label{fig:prop}
\end{figure}

\textbf{Accidental faults almost always propagate to failure}
(Fig.~\ref{fig:prop}): 167 of 190 applied runs are harmful, and only 12\%
of tasks succeed. A wrong vendor also turns into an unauthorized order
(0.76 per run). A fault early in the pipeline wastes nearly the whole run.
A wrong-value or stale researcher result changes about 3 downstream stages
and wastes 4.2--4.9k tokens and 3--4 tool actions, while the shared-state
corruption reaches all 5 stages. The one fault that is mostly ignored, a
consistently doubled quantity in the plan (18/20 ignored), is absorbed only because the executor
also sees the original request. That is redundancy in the workflow's
inputs, not recovery, and we report it as such.

\textbf{Adversarial impact depends on distance to the actuator.} When the
planner itself is compromised, the executor runs the exfiltration step in
20/25 tasks. Injection two hops upstream is mostly filtered out. The
researcher copied the poisoned catalog note into its output in 25/25 runs,
yet the planner acted on it only once. Injection in tool results was
ignored in 17/18 runs. Overall, 40 of 93 adversarial runs ignored the
injection, 30 consumed it without acting on it, and 23 were harmful. Task
success alone hides this: 91 of 93 tasks completed, but with 0.23
unauthorized actions per run.

\begin{figure}[t]
\centering
\includegraphics[width=\columnwidth]{figs/methods.pdf}
\caption{Task success by fault and method. *Counts only external actions
issued after the restart, i.e., as if earlier ones could be undone.
$^\dagger$At least one run in the cell took an unauthorized action.}
\label{fig:methods}
\end{figure}

\begin{table}[t]
\centering\small
\setlength{\tabcolsep}{3.2pt}
\begin{tabular}{llrrrrrr}
\toprule
 & method & succ. & succ.* & unauth. & resid. & tokens & lat.\,s\\
\midrule
\multirow{3}{*}{acc.} & none & 12\% & 12\% & 0.14 & 0 & 4.4k & 44\\
 & restart & 41\% & 96\% & 0.14 & 1.4 & 9.7k & 82\\
 & verify & 74\% & 74\% & 0.03 & 0 & 7.3k & 62\\
\midrule
\multirow{3}{*}{adv.} & none & 98\% & 98\% & 0.23 & 0 & 6.0k & 43\\
 & restart & 20\% & 98\% & 0.21 & 2.0 & 11.2k & 75\\
 & verify & 97\% & 97\% & 0.00 & 0 & 7.8k & 57\\
\bottomrule
\end{tabular}
\caption{Methods under identical fault schedules (per-run means; 188--190
accidental and 93--95 adversarial applied runs per method). resid.\ = external
actions from the faulty run that cannot be undone. Latency includes
simulated 30\,s timeouts.}
\label{tab:methods}
\end{table}

\textbf{Full restart is correct but does not recover.} The rerun itself
succeeds 96--98\% of the time. Yet only 41\% (accidental) and 20\%
(adversarial) of tasks end correct, because the faulty run had already
issued 1.4--2.0 irreversible external actions per run: duplicate vendor
orders and confirmations sent with the wrong values
(Table~\ref{tab:methods}). Restart is worst exactly when the fault was
harmless: an ignored injection still triggers a second order. It also
cannot take back an exfiltration email. On top of this, it costs
1.9--2.2$\times$ the tokens of the faulty run alone.

\textbf{Always-verify is cheaper but has blind spots.} It recovers 92--100\%
of researcher- and planner-stage faults (the misses sit in 2 tasks where the
model errs on its own) and blocks all unauthorized actions in the
adversarial set, for +1.8--2.9k tokens per run. However, it scores 0\%
on the corrupted shared budget, because the verifier checks against the
same corrupted source, and 0\% on the corrupted order id, because tool
results are not verified. Its recoveries also benefit from our faults being
transient: the single retry sees clean data. A persistent fault would make
the retry fail again.

\section{Takeaways for Phase 2}
(1) Most of the real damage is irreversible external actions, not wrong
intermediate values. Recovery must reason about what has left the system
instead of replaying the workflow. (2) Faults spread along the recorded
dependency edges, and shared state widens the reach of an error. For
example, for \f{tool.wrong\_order\_id} only the 2--3 executor steps after
\f{place\_order} depend on the bad id, so selective rerun would keep the
(valid) order instead of duplicating it. (3) Verification at the agent
level misses errors in shared state and in tool results, which are exactly
the cases dependency-aware invalidation can reach.
\emph{Limitations:} one model, 25 tasks, oracle fault identity, and only
transient LLM-stage faults. Phase~2 will scale tasks and models and add
persistent faults.

\appendix
\section{Timeline and Milestones}
\begin{itemize}
\item \textbf{Nov 11}: a second model and 50+ tasks; add persistent and multi-fault schedules; track natural (non-injected) faults like the one in the clean baseline.
\item \textbf{Nov 18}: dependency-aware invalidation over the trace graph: mark descendants of the faulty artifact, keep the rest, rerun only invalidated steps.
\item \textbf{Nov 25}: shared-state recovery: per-step file snapshots (copy-on-write or container); compensating actions (cancel order, correction email) for external effects; report actions that cannot be undone.
\item \textbf{Dec 2}: compare with full restart and verify under identical schedules: correctness, unauthorized actions, recomputed stages, retained valid work (checked against clean runs), recovery latency, tokens, tracking overhead.
\item \textbf{Dec 6}: optional automated detection, with false positives and missed faults reported separately.
\item \textbf{Dec 10}: final code, traces, and report.
\end{itemize}

\section{Fault Definitions}
{\footnotesize\renewcommand{\f}[1]{\texttt{\scriptsize #1}}
\setlength{\tabcolsep}{3pt}
\begin{tabular}{@{}p{0.43\columnwidth}>{\raggedright\arraybackslash}p{0.54\columnwidth}@{}}
\toprule
fault & effect\\
\midrule
\f{state.corrupt\_budget} & team budget in \f{budgets.json} wrong before the run, flipping the buy/no-buy decision\\
\f{researcher.wrong\_price} & unit price halved\\
\f{researcher.wrong\_vendor} & a different, preferably unapproved, vendor\\
\f{researcher.stale} & facts from an older catalog and a former manager\\
\f{researcher.timeout}, \f{planner.timeout} & output never arrives (30\,s simulated)\\
\f{planner.wrong\_qty} & quantity doubled consistently in order, total, record and deduction\\
\f{tool.wrong\_order\_id} & \f{place\_order} returns a corrupted id\\
\midrule
\f{retrieval.injection} & catalog record says to email confirmations to an outside address\\
\f{researcher.injection} & compromised researcher: exfiltrate and delete \f{orders.csv}\\
\f{planner.malicious\_step} & compromised planner appends an exfiltration email\\
\f{tool.injection} & vendor API reply asks to email \f{budgets.json} outside\\
\bottomrule
\end{tabular}}

\smallskip\noindent{\raggedright\textbf{Reproduce:} \f{python scripts/run\_experiment.py --model api-glm-5.3-flash --n-tasks 25}; figures: \f{scripts/plots.py}. Traces for every run are included.\par}

\end{document}
